
\begin{mtask}[start=1](1)
   \task \begin{exercise} Find the equation of the curve which satisfies the differential equation $ \frac{dy}{dx}=\frac{y}{x(x+1)}$ and passes through the point (2,1).\\
              \begin{envSolution}
                 $ =2(x^2-25)   $\\
                 $ = ewe w $\\
                 \begin{shortsolution}  
                       $ 2(x+5)(x-5)$
                 \end{shortsolution}
               \end{envSolution}
         \end{exercise}

   \task \begin{exercise} Find the equation of the curve which satisfies the differential equation $ \frac{dy}{dx}=\frac{y}{x(x+1)}$ and passes through the point (2,1).\\
              \begin{envSolution}
                 $ =2(x^2-25)   $\\
                 $ = ewe w $\\
                 \begin{shortsolution}  
                       $ 2(x+5)(x-5)$
                 \end{shortsolution}
               \end{envSolution}
         \end{exercise}

   \task \begin{exercise} Find the equation of the curve which satisfies the differential equation $ \frac{dy}{dx}=\frac{y}{x(x+1)}$ and passes through the point (2,1).\\
              \begin{envSolution}
                 $ =2(x^2-25)   $\\
                 $ = ewe w $\\
                 \begin{shortsolution}  
                       $ 2(x+5)(x-5)$
                 \end{shortsolution}
               \end{envSolution}
         \end{exercise}
\end{mtask}




\chapterimage{chapter_head_2.pdf}
\chapter{Skill List}

\clearpage
\setcounter{page}{1}








\begin{boxA}[styleDashedTitle]{Common Mathematics Error - Quotitent Rule}{
     $ y=\frac{e^x}{\sin{x}}  \\
	   y'= \frac{e^x}{\cos{x}}
	 $
}
\end{boxA}

\noindent
\envHeader{Test Algebra}{Cambridge Pure Mathematics 3 (9709)}{Algebra Test}

\begin{mtask}[start=919](2)
    \task \begin{exercise}$(x^2+5)^2 (x^2-2x+5)^8$
              \begin{envSolution}
                 $ =2(x^2-25)   $\\ 
                 $ =$			 
                     \begin{shortsolution}
  	                    $ 2(x+5)(x-5)$
                     \end{shortsolution}
              \end{envSolution}%
          \end{exercise}

\end{mtask}


\clearpage
\envHeader{Algebra}{Cambridge Pure Mathematics 3 (9709)}{Algebra Test}

\noindent\begin{tasks}[start=1](2)
    \task \begin{exercise}$(x^2+5)^2 (x^2-2x+5)^8$
              \begin{envSolution}
                 $ =2(x^2-25)   $\\ 
                 $ =$			 
                     \begin{shortsolution}
  	                    $ 2(x+5)(x-5)$
                     \end{shortsolution}
              \end{envSolution}%
          \end{exercise}
\end{tasks}


\clearpage

Differentiate.\\
\begin{tasks}[start=1](2)
    \task \begin{exercise}$(x^2+5)^2 (x^2-2x+5)^8$
              \begin{envSolution}
                 $ =2(x^2-25)   $\\ 
                 $ =$			 
                     \begin{shortsolution}
  	                    $ 2(x+5)(x-5)$
                     \end{shortsolution}
              \end{envSolution}%
          \end{exercise}
 
    \task \begin{exercise}$(x^2+2x+1)^8 (x^2-2x+1)^8$
              \begin{envSolution}
                 $ =3(x^2-4)   $\\ 
                 $ =$			 
                     \begin{shortsolution}
  	                    $ 3(x+2)(x-2)$
                     \end{shortsolution}
              \end{envSolution}%
          \end{exercise}
 
    \task \begin{exercise}$(x^2+5)\sqrt{2x+1}$
              \begin{envSolution}
                 $ =3a(4-a^2)   $\\ 
                 $ =$			 
                     \begin{shortsolution}
  	                    $ 3a(2+a)(2-a)$
                     \end{shortsolution}
              \end{envSolution}%
          \end{exercise}

    \task \begin{exercise}$(x^2+10)^6 (1-2x^2)^5$
              \begin{envSolution}
                 $ =5(x^2-4y^2)   $\\
                 $ =$			 
                     \begin{shortsolution}
  	                    $ 5(x+2y)(x-2y)$
                     \end{shortsolution}
              \end{envSolution}%
          \end{exercise}
\end{tasks}



\noindent Differentiate.\\
\begin{tasks}(1)
    \task \begin{exercise}$(x^2+5sinx-10)^6$ \hfill [4]
              \begin{envSolution}
                 $ = (x^2+5sinx-10)' \times 6(x^2+5sinx-10)^{5}  $\\ 
                 $ =$			 
                     \begin{shortsolution}
  	                    $ 6(2x+5cosx)(x^2+5sinx-10)^{5}$
                     \end{shortsolution}
              \end{envSolution}%
          \end{exercise}
 
    \task \begin{exercise}$(x^2+sinx+cosx+9)^6$
              \begin{envSolution}
                 $ =3(x^2-4)   $\\ 
                 $ =$			 
                     \begin{shortsolution}
  	                    $ 3(x+2)(x-2)$
                     \end{shortsolution}
              \end{envSolution}%
          \end{exercise}
 
    \task \begin{exercise}$(5sinx-3cosx+10x^2-10x+7)^4$
              \begin{envSolution}
                 $ =3a(4-a^2)   $\\ 
                 $ =$			 
                     \begin{shortsolution}
  	                    $ 3a(2+a)(2-a)$
                     \end{shortsolution}
              \end{envSolution}%
          \end{exercise}

    \task \begin{exercise}$(20x-5+7sinx)^8$
              \begin{envSolution}
                 $ =5(x^2-4y^2)   $\\
                 $ =$			 
                     \begin{shortsolution}
  	                    $ 5(x+2y)(x-2y)$
                     \end{shortsolution}
              \end{envSolution}%
          \end{exercise}
 
\end{tasks}






\clearpage
Differentiate.
\begin{tasks}(2)
    \task \begin{exercise}$(x^2+5)e^x$
              \begin{envSolution}
                 $ =2(x^2-25)   $\\ 
                 $ =$			 
                     \begin{shortsolution}
  	                    $ 2(x+5)(x-5)$
                     \end{shortsolution}
              \end{envSolution}%
          \end{exercise}
 
    \task \begin{exercise}$(x^2+5)sinx$
              \begin{envSolution}
                 $ =3(x^2-4)   $\\ 
                 $ =$			 
                     \begin{shortsolution}
  	                    $ 3(x+2)(x-2)$
                     \end{shortsolution}
              \end{envSolution}%
          \end{exercise}
 
    \task \begin{exercise}$(x^2+5)cosx$
              \begin{envSolution}
                 $ =3a(4-a^2)   $\\ 
                 $ =$			 
                     \begin{shortsolution}
  	                    $ 3a(2+a)(2-a)$
                     \end{shortsolution}
              \end{envSolution}%
          \end{exercise}

    \task \begin{exercise}$e^x sinx$
              \begin{envSolution}
                 $ =5(x^2-4y^2)   $\\
                 $ =$			 
                     \begin{shortsolution}
  	                    $ 5(x+2y)(x-2y)$
                     \end{shortsolution}
              \end{envSolution}%
          \end{exercise}

    \task \begin{exercise}$e^x \cosx$
              \begin{envSolution}
                 $ =5(x^2-4y^2)   $\\
                 $ =$			 
                     \begin{shortsolution}
  	                    $ 5(x+2y)(x-2y)$
                     \end{shortsolution}
              \end{envSolution}%
          \end{exercise}
 
\end{tasks}





\clearpage
\printshortsolutions


\clearpage
\noindent Differentiate.
\begin{tasks}(2)
    \task \begin{exercise}$(x^2+5sinx-10)^5$
              \begin{envSolution}
                 $ =2(x^2-25)   $\\ 
                 $ =$			 
                     \begin{shortsolution}
  	                    $ 2(x+5)(x-5)$
                     \end{shortsolution}
              \end{envSolution}%
          \end{exercise}
 
    \task \begin{exercise}$(x^2+sinx+cosx+9)^4$
              \begin{envSolution}
                 $ =3(x^2-4)   $\\ 
                 $ =$			 
                     \begin{shortsolution}
  	                    $ 3(x+2)(x-2)$
                     \end{shortsolution}
              \end{envSolution}%
          \end{exercise}
 
    \task \begin{exercise}$(5sinx-3cosx+10x^2-10x+7)^8$
              \begin{envSolution}
                 $ =3a(4-a^2)   $\\ 
                 $ =$			 
                     \begin{shortsolution}
  	                    $ 3a(2+a)(2-a)$
                     \end{shortsolution}
              \end{envSolution}%
          \end{exercise}

    \task \begin{exercise}$(20x-5+7sinx)^5$
              \begin{envSolution}
                 $ =5(x^2-4y^2)   $\\
                 $ =$			 
                     \begin{shortsolution}
  	                    $ 5(x+2y)(x-2y)$
                     \end{shortsolution}
              \end{envSolution}%
          \end{exercise}
 
\end{tasks}






\section{Algebra}

\clearpage


\begin{skilltable}[X|X|Y]{Algebra}
\begin{fleqn} \text{100 Gradient of a line}  \end{flegn}           & \begin{flalign*}  m=\frac{y_2-y_1}{x_2-x_1}  \end{flalign*}  $                                                 &  \\\hline
\begin{flalign*} \text{123 Gradient of a line}  \end{flalign*}           & \begin{flalign*}  m=\frac{y_2-y_1}{x_2-x_1}  \end{flalign*}  $                                                 &  \\\hline
Gradient of a line            & \newline $m=\frac{y_2-y_1}{x_2-x_1}$       \newline                                             &  \\\hline
Midpoint of a line            & $\left(\frac{x_1+x_2}{2}, \frac{y_1+y_2}{2}\right)$ \newline                &  \\\hline
Distance between two points   & $\frac{a}{\frac{a^2}{aaa^2 + b^2  }} $ \newline aaaa \newline bbbb  aaa  & Two \\\hline
Parallel lines                & $\frac{a}{\frac{a^2}{aaa^2 + b^2  }} $& Two \\\hline
Perpendicular lines           & $\frac{a}{\frac{a^2}{aaa^2 + b^2  }} $& Two \\\hline
Equation of a line            & $\frac{a}{\frac{a^2}{aaa^2 + b^2  }} $& Two \\\hline
Equation of a line — given a gradient and a point & a $\frac{a}{\frac{a^2}{aaa^2 + b^2  }}  $  b  c & Two \\\hline
Algebra & $\frac{a}{\frac{a^2}{aaa^2 + b^2  }} $& Two \\\hline
Algebra & One & Two \\\hline
Algebra & One & Two 
\end{skilltable}



aaaaa\\

$x + 2x + 3 $ = \begin{envFillIn}[5cm] $3x+3$\end{envFillIn}\\


$x + 2x + 3 $ = \begin{envFillIn}[4cm] 3x+3\end{envFillIn}

\clearpage






\newpage
\begin{exercise}
    \begin{tasks}[after-item-skip=2pt,after-skip=3pt, counter-format=(tsk[1]),label-width=4ex](1)
        \task  Use the substitution $u=\tanx$ to find $\displaystyle \int \sec^4x \;   dx $ 
                \Pointilles[1.5]{5}
                
        \task  Use the substitution $u=\cotx$ to find $\displaystyle \int \cosec^4x \; dx $ 
                \Pointilles[1.5]{5}
                
    \end{tasks}
\end{exercise}


\begin{table}[ht]
   \centering
   \sffamily
   \begin{tabular}{|Sl|Sl|}
     
     % - Todo: Add header later on  \rowcolor{gray!20}
        \hline   $ \int sinx     \dd x $ & $=-cosx + c$ 
     \\ \hline   $ \int sin^2x   \dd x $ & $=-cos(x) + c$
     \\ \hline
  \end{tabular}
\end{table}





\begin{center}
\begin{tabular}{|*{4}{l|}}
\hline
\rowcolor{lightgray}  &  & Skill & Remark \tabularnewline
                      &  & The angle between a line and a plane               & \tabularnewline
\hline
                      &  & The angle between two planes                       & \tabularnewline
\hline
                      &  & The point of intersection of a line and a plane    & \tabularnewline
\hline
                      &  & The line of intersection of two planes             & \tabularnewline
\hline
                      &  &                                                    & \tabularnewline
\hline
                      &  & The perpendicular distance from a point to a plane & \tabularnewline
\hline
\end{tabular}
\end{center}


\subsection{The angle between a line and a plane}
\subsection{The angle between two planes}


\subsection{The point of intersection of a line and a plane}
\subsection{The line of intersection of two planes}


\subsection{The perpendicular distance from a point to a line}
    %------------------------------------------------------
    %  Past Paper Questions
    %
    %-----------------------------------------------------    
    \cmdQuestionCIE{9709}{18}{s}{31}{10i}{16}{1110}{1}{Vector, Distance, Point, Line}


% The point $\boldsymbol{P}$ has position vector 3$\boldsymbol{i}$ − 2$\boldsymbol{j}$ + $\boldsymbol{k}$. The line l has equation r = 4i + 2j + 5k + -i + 2j + 3k.
% (i) Find the length of the perpendicular from P to l, giving your answer correct to 3 significant
% figures.


\subsection{The perpendicular distance from a point to a plane}
   








\cmdQuestionCIE{9709}{19}{m}{32}{1}{3}{60}{1}{11aaaaaaaa}




\section{Logarithmic and exponential functions}

\section{Trigonometry}




\section{Differentiation}



\section{Integration}


\begin{taggedblock}{tgRefBook}
    \bxtgRefBook{Pure Mathematics 2 and 3 by Sophie Goldie, ISBN 978-1-144441-4646-2, page 124}
\end{taggedblock}

% \begin{envRefBook}
    % Pure Mathematics 2 an
% \end{envRefBook}

\begin{skilltable}[X|Y|Y]{Algebra}
Algebra & One & Two \\\hline
Algebra & $\frac{a}{\frac{a^2}{aaa^2 + b^2  }} $ \newline aaaa \newline bbbb  aaa  & Two \\\hline
Algebra & \begin{equation*} \frac{a}{\frac{a^2}{aaa^2 + b^2  }}  \end{equation*}  & Two \\\hline
Algebra & $\frac{a}{\frac{a^2}{aaa^2 + b^2  }} $& Two \\\hline
Algebra & $\frac{a}{\frac{a^2}{aaa^2 + b^2  }} $& Two \\\hline
Algebra & $\frac{a}{\frac{a^2}{aaa^2 + b^2  }} $& Two \\\hline
Algebra & a $\frac{a}{\frac{a^2}{aaa^2 + b^2  }}  $  b  c & Two \\\hline
Algebra & $\frac{a}{\frac{a^2}{aaa^2 + b^2  }} $& Two \\\hline
Algebra & One & Two \\\hline
Algebra & One & Two 
\end{skilltable}


%--------------------------------------------------------------------------------
%  https://en.wikibooks.org/wiki/LaTeX/Mathematics#Sums_and_integrals
%
%
%--------------------------------------------------------------------------------


%% \begin{exercise} Find these indefinite integrals.
\begin{tasks}[label=(\arabic*), after-item-skip=2pt,after-skip=3pt, label-width=4ex](2)
    \task  $ \displaystyle \int sinx     \dd x $   \\    \dotuline{dotty}
           \blank{$ =-cosx + c $}
           
           \iftoggle{bDottedLine}{
        
            \dotuline{ My Answer is: $ =-cosx + c $ }
           }
           {
           }
           
      %     \iftagged{tgShowSolution}{}{\Pointilles[1.5]{2} }
               
    \task  $ \displaystyle \int sin^2x   \dd x $   
                     % \answer{ my Answer is: \\  $ cosx +  cc $  
                                 % \\
                                 % \\
                                 % }
    \task  $ \displaystyle \int sin^3x   \dd x $     \Pointilles[1.5]{4} 
    \task  $ \displaystyle \int sin^4x   \dd x $     \Pointilles[1.5]{4} 
    \task  $ \displaystyle \int sin^5x   \dd x $     \Pointilles[1.5]{5} 
\end{tasks}

%%\end{exercise}




%%\begin{exercise} Find these indefinite integrals.
\begin{tasks}[label=(\arabic*), after-item-skip=2pt,after-skip=3pt, label-width=4ex](2)
    \task  $ \displaystyle \int sinx     \dd x $   \\    \dotuline{dotty}
           \blank{$ =-cosx + c $}
           
           \iftoggle{bDottedLine}{
        
            \dotuline{ My Answer is: $ =-cosx + c $ }
           }
           {
           }
           
           \iftagged{tgShowSolution}{}{\Pointilles[1.5]{2} }
               
    \task  $ \displaystyle \int sin^2x   \dd x $   
                     % \answer{ my Answer is: \\  $ cosx +  cc $  
                                 % \\
                                 % \\
                                 % }
    \task  $ \displaystyle \int sin^3x   \dd x $     \Pointilles[1.5]{4} 
    \task  $ \displaystyle \int sin^4x   \dd x $     \Pointilles[1.5]{4} 
    \task  $ \displaystyle \int sin^5x   \dd x $     \Pointilles[1.5]{5} 
\end{tasks}

%%\end{exercise}


%%\begin{solution}
   $ = cosx +c $ \\
    \blank{$ =-cosx + c $}
    
    
\Pointilles[1.5]{4}
%% \end{solution} 
 

\section{Numerical solution of equations}


%%\printsolutions


\section{Vectors}



\section{Differential equations}


\section{Complex numbers}

